\documentclass{poly}
\usepackage{main}

\title{Correction}
\author{Seconde 3}
\date{16 Février 2026}

\begin{document}
\maketitle

On commence par faire une figure de la situation exposée dans l'énoncé.

\begin{center}
\begin{tikzpicture}
\draw[help lines] (-3.25,-3.25) grid[step=0.5] (3.25,3.25);
\coordinate (A) at (-2, 2);
\coordinate (B) at (3, 2);
\coordinate (C) at (0,-2);
\coordinate (D) at (-1,-2);
\coordinate (I) at (0.5,2);
\coordinate (K) at (-0.5,-2);
\draw (A) -- (B) -- (C) -- (D) -- cycle;
\draw (A) node[above left] {$A$};
\draw (B) node[above right] {$B$};
\draw (C) node[below right] {$C$};
\draw (D) node[below left] {$D$};
\draw[name path = AC] (A) -- (C);
\draw[name path = DB] (D) -- (B);
\draw (I) node {$\bullet$} node[above] {$I$};
\draw (K) node {$\bullet$} node[below] {$K$};
\draw[dashed] (I) -- (K);
\path [name intersections={of=AC and DB, by=M}];
\draw (M) node {$\bullet$} node[right] {$M$};
\end{tikzpicture}
\end{center}

On conjecture que les points $K$, $M$ et $I$ sont alignés.

On prouve cette conjecture à l'aide des vecteurs. On cherche donc à montrer que les vecteurs $\vect{KM}$ et $\vect{IM}$ sont colinéaires.

On traduit ce que l'on sait grâce à l'énoncé :
\begin{itemize}
\item On sait déjà que $\vect{AB} = 5\vect{DC}$;
\item $I$ est le milieu de $[AB]$ : on sait donc d'après le cours que $\vect{AI} = \vect{IB}$, ou que $\vect{AI}=\dfrac{1}{2}\vect{AB}$.
\item $K$ est le milieu de $[DC]$ : on sait donc d'après le cours que $\vect{DK} = \vect{KC}$, ou que $\vect{DK}=\dfrac{1}{2}\vect{DC}$.
\item $M$ est le point d'intersection de $[AC]$ et $[DB]$ : plus difficile, on peut simplement dire d'un côté que $A$, $M$ et $C$, et de l'autre côté que $D$, $M$ et $B$ sont alignés. Du point des vecteurs, on en déduit que $\vect{AM}$ et $\vect{MC}$ sont colinéaires, et que $\vect{BM}$ et $\vect{MD}$ sont colinéaires.
\end{itemize}

% On choisit de partir de $\vect{IM}$, et de trouver que ce vecteur est égal à $k\vect{KM}$ où $k$ sera un certain nombre. Le but est de faire apparaître les vecteurs sur lesquel on a des informations. Le vecteur $\vect{IM}$ peut-être obtenu par relation de Chasles de deux manières différentes. 

% \begin{equation*}
% \begin{cases}
% \vect{IM} &= \vect{IA} + \vect{AM}\\
% \vect{IM}&= \vect{IB} + \vect{BM}\\
% \end{cases}
% \end{equation*}
% On en déduit, en ajoutant ces deux lignes, que
% \begin{equation*}
% \begin{aligned}
% 2\vect{IM} &= \vect{IA} + \vect{AM} + \vect{IB} + \vect{BM}\\
% &= \vect{IA} + \vect{AI} + \vect{AM} + \vect{BM}\\
% &= \vect{0} + \vect{AM} + \vect{BM}\\
% &= \vect{AM} + \vect{BM}\\
% \end{aligned}
% \end{equation*}

\begin{equation*}
\begin{aligned}
\vect{MI} &= \vect{MA} + \vect{AI}\\
&= \vect{MA} + \dfrac{1}{2}\vect{AB}\\
&= k\vect{CM} + \dfrac{1}{2}(5\vect{DC})\\
&= k\vect{CM} + 5(\dfrac{1}{2}\vect{DC})\\
&= k\vect{CM} + 5\vect{KC}
\end{aligned}
\end{equation*}

Si $k = 5$, alors on peut factoriser par $5$ :
\begin{equation*}
\vect{MI} = 5\vect{CM} + 5\vect{KC} = 5(\vect{KC} + \vect{CM}) = 5\vect{KM}  
\end{equation*}

Cela nous permet de montrer que $\vect{MI}$ et $\vect{KM}$ sont colinéaires, ce qui correspond à ce que nous cherchons.

Mais cela ne fonctionne que si $k = 5$, il faut donc le prouver. 

Le $5$ nous pousse à prendre comme point de départ l'égalité $\vect{AB} = 5\vect{DC}$.

\begin{equation*}
\begin{aligned}
&\vect{AB} &= 5\vect{DC}\\
\Leftrightarrow&\vect{AM} + \vect{MB} &= 5\vect{DM} + 5\vect{MC}\\
\Leftrightarrow&k\vect{MC} + k'\vect{DM} &= 5\vect{DM} + 5\vect{MC}\\
\Leftrightarrow&(k-5)\vect{MC} &= (5 - k')\vect{DM}\\
\end{aligned}
\end{equation*}
On sait que $\vect{MC}$ et $\vect{DM}$ ne sont pas colinéaires. Donc cette égalité qui prétend le contraire ne peut être vraie que si $(k - 5) = (5 - k') = 0$. Donc $k = k' = 5$. Ainsi, on en déduit exactement ce que l'on voulait prouver.


\end{document}